\section{The \arb{} Benchmark}\label{sec:bench}

\subsection{A compact, signal-defined failure taxonomy}

\arb{} uses exactly seven failure classes. The design constraint was that a class must be
(1) \emph{observable}---there exists a monitor-visible signal that discriminates it; (2)
\emph{injectable}---we can realise it deterministically inside a tool environment; and (3)
\emph{actionable}---a specific repair operator is prescribed. Table~\ref{tab:taxonomy} lists them
with detectability, canonical remediation, and recoverability.

\begin{table*}[t]
\centering
\small
\caption{The \arb{} failure taxonomy. \emph{Detectability} is a property of the signal, not of a
model: \emph{surface} classes are visible in error codes or version counters, \emph{latent} classes
require cross-referencing evidence or the objective specification. \emph{Canonical remediation} is
the operator a domain expert prescribes; it defines the rule-based reference policy.}
\label{tab:taxonomy}
\begin{tabular}{@{}lllll@{}}
\toprule
Class & Detectability & Track & Canonical remediation & Recoverable \\
\midrule
invalid\_arguments & surface & single+multi & \texttt{revalidate\_args} & yes \\
stale\_state & surface & single+multi & \texttt{re\_read\_authoritative} & yes \\
contradictory\_evidence & latent & single+multi & \texttt{refresh\_state} & yes \\
planning\_error & latent & single+multi & \texttt{replan} & yes \\
communication\_failure & surface & multi & \texttt{resync\_message} & yes \\
objective\_drift & latent & single+multi & \texttt{reanchor\_objective} & yes \\
unrecoverable\_external\_error & surface & single+multi & \texttt{escalate\_abort} & no (contain) \\
\bottomrule
\end{tabular}
\end{table*}

The split between \emph{surface} and \emph{latent} is the load-bearing distinction. A surface class
announces itself: the tool returns a schema violation, or a read carries a version counter older
than the current one. A latent class does not: contradictory evidence is two {\em individually
plausible\/} values for the same field, planning error is a well-formed step that serves the wrong
subgoal, and objective drift is a locally sensible action that no longer serves the stated goal.
Latent classes require an internal model of what the task \emph{should} be doing, which is precisely
what a monitor built only on error codes lacks.

\texttt{escalate\_abort} is deliberately not a repair. For an unrecoverable external error the
correct behaviour is to stop mutating, leave the world consistent, and hand off. We therefore score
containment (no successful mutation after the fault) separately from recovery, and we never count a
containment as a recovery.

\subsection{Environments and the reference agent}

Four families, each a deterministic Python world with a version counter and a bounded history:

\begin{itemize}
\item \textbf{Refund desk (\texttt{tickets}).} Support tickets with a refund policy whose arithmetic
(constituent caps, prior refunds, order totals) must be recomputed each time. The SOP includes a
check-then-act re-read before moving money, which makes stale-read faults physically costly.
\item \textbf{Dual-store migration (\texttt{migration}).} A source store and a migrated store; the
agent copies records and must reconcile. Background writers touch only objects the agent has not
read, so silent lost updates are possible and a completion claim can be wrong.
\item \textbf{Inventory accounting (\texttt{inventory}).} An auditable invariant
(\texttt{on\_hand} $+$ $\sum$ \texttt{adjustments} $=$ \texttt{ledger}) plus a \texttt{close\_period}
gate. Corrections are made by posting compensating adjustments, never by editing history.
\item \textbf{Procurement desk (\texttt{procurement}).} Three roles---planner, buyer, verifier---on a
message bus with an incremental order protocol (cancel and re-issue is the repair path). This is the
multi-agent family; the others are single-agent.
\end{itemize}

The \emph{reference agent} is a scripted workflow executor. It knows the procedure and the
environment's affordances (\texttt{solve}/\texttt{next\_step}) but has \emph{no access to ground
truth}: every tool argument is derived from that role's belief, i.e.\ from observations it has
actually received and applied. This design choice is what makes the hard classes real. If the
argument were computed from the world, a stale read would be harmless; because it is computed from
belief, a stale read propagates into the next mutation. Similarly, a dropped message is not merely
``an event we log''---the receiving role's belief never updates, so it acts on the wrong sender or
the wrong quantity.

Agent-side ``natural'' faults (omitting a required verification step, using a superseded value,
mis-ordering a two-phase mutation) can be switched on to produce failures that were not injected by
the harness; we use them as a provenance-shift check (Table~\ref{tab:natural}).

\subsection{Tasks, splits, and leakage control}

A task is a (family, template, configuration, seed) tuple. Templates define the procedure shape;
configurations vary concurrency, dataset size, and which optional branches fire. Splits hold out
both:

\begin{itemize}
\item \textbf{train} --- templates A/B $\times$ configs 1--3;
\item \textbf{dev} --- templates A/B $\times$ config 5;
\item \textbf{test-id} --- templates A/B $\times$ config 6 (unseen configuration, higher concurrency);
\item \textbf{test-ood} --- template C (unseen template).
\end{itemize}

Detector thresholds are tuned on \emph{dev}; the supervised detector and the recovery ranker are fit
on \emph{train} only. Nothing in test-id or test-ood is seen during fitting, and fault schedules are
generated per task from a seeded RNG so that memorising a schedule is not a viable strategy.

\subsection{Fault injection as a program}

Each class is realised by a small \emph{operator program} applied at a scheduled step: a sequence of
primitive perturbations (corrupt an argument, re-serve an older snapshot, plant a conflicting value,
suppress a message, drop a subgoal, swap the objective, raise a terminal tool error). \texttt{inject\_at}
is aligned to the step that calls the targeted tool---the last read of that field, or the first
mutation---so a scheduled fault is always actually exercised. The same program is used for every
family, which is what lets us pool classes across environments.

\subsection{Harm screening}

Not every injection changes the outcome. A stale read that is never used, or a duplicated message
that is idempotent, leaves the task successful. We therefore run every (task, fault) pair under
\texttt{no\_recovery} and keep only pairs where the task actually fails: \NumHarmful{} of
\NumPairs{} (\InertPct\% inert, discarded). This has two effects. It defines the denominator, so
``recovery rate'' means ``fraction of \emph{harmful} faults repaired''. And it is conservative for
detection: faults that are present but harmless are never counted as false negatives or as
detections. The cost is a selection effect---the retained set is biased toward faults with large
downstream reach---which we discuss in Sec.~\ref{sec:limits}.
